%% kapa: exam=1 round=20 year=2009 period=1 form=A subject=civil_law subject_ko=민법 seq=6 number=6
%% source: Q-Net 감정평가사 기출문제 (cst00309) artlSeq=5024178 file=extracted/1008/5024178/2009년 제20회 감정평가사 1차시험 기출문제/1교시.hwp
%% conversion: kordoc 4.24.51 -> tools/gen1.py (auto); review=proofread-sonnet; flags=none
\begin{question}{6}
공동저당에 관한 설명으로 옳지 \uline{않은} 것은? (다툼이 있으면 판례에 의함)
\begin{choices}
  \item 대지와 건물이 동시에 경매되어 공동저당권자에게 그 경매대가를 동시에 배당하는 때에는, 대지와 건물의 경매대가에 비례하여 그 채권의 분담을 정하여야 한다.
  \item ①에서 부담안분의 원칙은 저당부동산에 대하여 후순위저당권자가 존재하는가의 여부를 불문하고 적용된다.
  \item 채무자의 부동산보다 물상보증인의 부동산에 대하여 먼저 담보권이 실행되었다면, 물상보증인의 부동산에 후순위저당권을 가진 자보다 물상보증인이 우선적으로 보호된다.
  \item 공동저당권자는 임의로 어느 저당목적물로부터 채권의 전부나 일부의 우선변제를 받을 수 있다.
  \item 공동저당부동산이 5개 이상일 때에는 공동담보목록을 첨부하여야 한다.
\end{choices}
\end{question}
\answer{3}
